% GENERATED FILE. Edit canonical lessons, then run npm run generate:books.
% canonical-source-hash: fnv1a64:c02c63cd22329841
% canonical-lessons: FR-C452-remise, FR-C452-rabais, FR-C452-marchander, FR-C452-brader, FR-C452-braderie

\chapter{Les bonnes affaires}
\label{ch:fr-les-bonnes-affaires}
\hlchaptermodality{452}

\begin{chapteropening}
\textbf{By the end of this chapter:} \emph{I can find a bargain: the discount, the price cut, haggling, selling off and the street sale.}

Complete the last lesson of chapter 452: I can find a bargain: the discount, the price cut, haggling, selling off and the street sale.
\end{chapteropening}

\section[la remise]{la remise — the discount; the shed}

\label{lesson:FR-C452-remise}



\begin{quote}
Before the new one: say the French for the overdraft, then the French for the statement.
\end{quote}

\subsection*{You'll want to know: la remise}
\textbf{la remise} — \textquotedblleft{}the discount; the shed\textquotedblright{}.

\begin{quote}
Le vendeur m'a fait une remise de dix pour cent.
\end{quote}

\begin{quote}
\emph{The salesman gave me ten per cent off.}
\end{quote}

\emph{Une remise} is also a garden shed.

\begin{cousinweb}
\emph{remise} is the feminine past participle of \emph{remettre} (EtymDB), from Latin \emph{remittere} 'to send back', built on \emph{mittere} 'to send'.
\end{cousinweb}

\begin{grammarlens}[title={What you've built}]
One more word for telling, explaining and describing at B1.
\end{grammarlens}

\subsection*{Your turn}
\begin{itemize}
\raggedright
  \item \emph{Say it:} \emph{la remise}
  \item \emph{Say it:} \emph{la remise}, once more
  \item \emph{Say it:} a sentence of your own with \emph{la remise}
  \item \emph{Recall:} say the French for the overdraft, then the French for the statement, then say \emph{la remise} again
\end{itemize}

\begin{tcolorbox}[breakable,colback=teal!4,colframe=teal!35!black,title={Before you move on}]
What does la remise mean? (\textquotedblleft{}The discount; the shed\textquotedblright{}.) Say it once more.
\end{tcolorbox}
\section[le rabais]{le rabais — the reduction, the price cut}

\label{lesson:FR-C452-rabais}



\begin{quote}
Before the new one: say the French for the statement, then the French for the discount.
\end{quote}

\subsection*{You'll want to know: le rabais}
\textbf{le rabais} — \textquotedblleft{}the reduction, the price cut\textquotedblright{}.

\begin{quote}
Ces chaussures sont vendues au rabais.
\end{quote}

\begin{quote}
\emph{These shoes are being sold at a reduced price.}
\end{quote}

\emph{Au rabais} — cut-price, sometimes cheap and poor.

\begin{cousinweb}
\emph{rabais} is built on \emph{rabaisser} 'to bring down', from \emph{baisser} 'to lower', which goes back to Vulgar Latin \emph{bassiāre}, from Late Latin \emph{bassus} 'low'. EtymDB does not list it.
\end{cousinweb}

\begin{grammarlens}[title={What you've built}]
One more word for telling, explaining and describing at B1.
\end{grammarlens}

\subsection*{Your turn}
\begin{itemize}
\raggedright
  \item \emph{Say it:} \emph{le rabais}
  \item \emph{Say it:} \emph{le rabais}, once more
  \item \emph{Say it:} a sentence of your own with \emph{le rabais}
  \item \emph{Recall:} say the French for the statement, then the French for the discount, then say \emph{le rabais} again
\end{itemize}

\begin{tcolorbox}[breakable,colback=teal!4,colframe=teal!35!black,title={Before you move on}]
What does le rabais mean? (\textquotedblleft{}The reduction, the price cut\textquotedblright{}.) Say it once more.
\end{tcolorbox}
\section[marchander]{marchander — to haggle, to bargain}

\label{lesson:FR-C452-marchander}



\begin{quote}
Before the new one: say the French for the discount, then the French for the reduction.
\end{quote}

\subsection*{You'll want to know: marchander}
\textbf{marchander} — \textquotedblleft{}to haggle, to bargain\textquotedblright{}.

\begin{quote}
Au marché aux puces, tout le monde marchande.
\end{quote}

\begin{quote}
\emph{At the flea market everyone haggles.}
\end{quote}

\emph{Le marchandage} — haggling.

\begin{cousinweb}
\emph{marchander} is built on \emph{marchand} 'trader' (EtymDB), which goes back to a Vulgar Latin form of \emph{mercārī} 'to trade', from \emph{merx} 'goods'.
\end{cousinweb}

\begin{grammarlens}[title={What you've built}]
One more word for telling, explaining and describing at B1.
\end{grammarlens}

\subsection*{Your turn}
\begin{itemize}
\raggedright
  \item \emph{Say it:} \emph{marchander}
  \item \emph{Say it:} \emph{marchander}, once more
  \item \emph{Say it:} a sentence of your own with \emph{marchander}
  \item \emph{Recall:} say the French for the discount, then the French for the reduction, then say \emph{marchander} again
\end{itemize}

\begin{tcolorbox}[breakable,colback=teal!4,colframe=teal!35!black,title={Before you move on}]
What does marchander mean? (\textquotedblleft{}To haggle, to bargain\textquotedblright{}.) Say it once more.
\end{tcolorbox}
\section[brader]{brader — to sell off cheap}

\label{lesson:FR-C452-brader}



\begin{quote}
Before the new one: say the French for the reduction, then the French for to haggle.
\end{quote}

\subsection*{You'll want to know: brader}
\textbf{brader} — \textquotedblleft{}to sell off cheap\textquotedblright{}.

\begin{quote}
En janvier, les boutiques bradent leurs vêtements d'hiver.
\end{quote}

\begin{quote}
\emph{In January the shops sell off their winter clothes cheap.}
\end{quote}

Also figurative: \emph{brader son talent} — to throw it away.

\begin{cousinweb}
\emph{brader} is usually traced to Dutch \emph{braden} 'to roast', from the meat roasted at old street markets; EtymDB lists it without a source, and the details are not settled.
\end{cousinweb}

\begin{grammarlens}[title={What you've built}]
One more word for telling, explaining and describing at B1.
\end{grammarlens}

\subsection*{Your turn}
\begin{itemize}
\raggedright
  \item \emph{Say it:} \emph{brader}
  \item \emph{Say it:} \emph{brader}, once more
  \item \emph{Say it:} a sentence of your own with \emph{brader}
  \item \emph{Recall:} say the French for the reduction, then the French for to haggle, then say \emph{brader} again
\end{itemize}

\begin{tcolorbox}[breakable,colback=teal!4,colframe=teal!35!black,title={Before you move on}]
What does brader mean? (\textquotedblleft{}To sell off cheap\textquotedblright{}.) Say it once more.
\end{tcolorbox}
\section[la braderie]{la braderie — the street sale, the clearance market}

\label{lesson:FR-C452-braderie}



\begin{quote}
Before the new one: say the French for to haggle, then the French for to sell off cheap.
\end{quote}

\subsection*{You'll want to know: la braderie}
\textbf{la braderie} — \textquotedblleft{}the street sale, the clearance market\textquotedblright{}.

\begin{quote}
Chaque septembre, la ville organise une grande braderie.
\end{quote}

\begin{quote}
\emph{Every September the town holds a big street sale.}
\end{quote}

Very popular in the north of France and in Belgium.

\begin{cousinweb}
\emph{braderie} is built on \emph{brader} 'to sell off cheap' (EtymDB), which is usually traced to Dutch \emph{braden} 'to roast'.
\end{cousinweb}

\begin{grammarlens}[title={What you've built}]
One more word for telling, explaining and describing at B1.
\end{grammarlens}

\subsection*{Your turn}
\begin{itemize}
\raggedright
  \item \emph{Say it:} \emph{la braderie}
  \item \emph{Say it:} \emph{la braderie}, once more
  \item \emph{Say it:} a sentence of your own with \emph{la braderie}
  \item \emph{Recall:} say the French for to haggle, then the French for to sell off cheap, then say \emph{la braderie} again
\end{itemize}

\begin{tcolorbox}[breakable,colback=teal!4,colframe=teal!35!black,title={Before you move on}]
What does la braderie mean? (\textquotedblleft{}The street sale, the clearance market\textquotedblright{}.) Say it once more.
\end{tcolorbox}
